\begin{abstractind}

Metode \textit{Laser-Induced Breakdown Spectroscopy} (LIBS) memungkinkan analisis komposisi elemental material secara cepat dan tanpa preparasi sampel, menjadikannya kandidat yang tepat untuk karakterisasi geokimia tanah vulkanik. Namun, akurasi kuantifikasi konvensional seringkali terhambat oleh tumpang tindih garis emisi, efek matriks, dan fluktuasi plasma yang kompleks. Penelitian ini mengusulkan kerangka analisis berbasis \textit{Deep Learning} menggunakan arsitektur CNN--Transformer \textit{encoder}--\textit{decoder} dengan mekanisme \textit{Cross-Attention} untuk dekomposisi spektrum campuran menjadi komponen unsur tunggal per elemen. Spektrum sintetis LIBS dibangkitkan menggunakan model plasma berbasis fisika melalui persamaan Saha--Boltzmann dalam kerangka \textit{Local Thermodynamic Equilibrium} (LTE), dengan memanfaatkan parameter transisi atom dari \textit{Atomic Spectra Database} (ASD) NIST. Mekanisme \textit{Cross-Attention} menggunakan \textit{embedding} parameter garis atomik NIST sebagai \textit{query} untuk memandu proses dekomposisi spektral secara terarah. Untuk menjembatani kesenjangan distribusi antara domain spektrum sintetis dan spektrum LIBS lapangan dari tanah vulkanik Gunung Sinabung, diterapkan adaptasi domain berbasis \textit{Gradient Reversal Layer} (GRL). Performa model dievaluasi terhadap hasil pengukuran \textit{X-ray Fluorescence} (XRF) sebagai acuan, dengan target $R^2 > 0{,}95$ dan MAPE $< 10\%$ pada elemen pembentuk batuan utama. Pendekatan ini diharapkan memberikan solusi yang terukur dan mudah diinterpretasi untuk pemantauan geokimia vulkanik secara \textit{in-situ} di Indonesia.

\bigskip
\noindent
\textbf{Kata kunci :} LIBS, CNN-Transformer, Spektrum Sintetis, Adaptasi Domain, \textit{Cross-Attention}, Tanah Vulkanik.
\end{abstractind}

